\documentclass[12pt, titlepage]{article}

\usepackage{booktabs}
\usepackage{tabularx}
\usepackage{hyperref}
\hypersetup{
    colorlinks,
    citecolor=blue,
    filecolor=black,
    linkcolor=red,
    urlcolor=blue
}
\usepackage[round]{natbib}
\usepackage{longtable}
\usepackage{booktabs}

\input{../Comments.text}
\input{../Common.text}

% Zero-size, empty anchor placed before the term, so it can't shift the text
\providecommand{\glsentry}[2]{}
\renewcommand{\glsentry}[2]{\raisebox{0pt}[0pt][0pt]{\hypertarget{gls:#1}{}}#2}
\providecommand{\glink}[2]{\hyperlink{gls:#1}{#2}}

\begin{document}

\title{System Verification and Validation Plan for \progname{}} 
\author{\authname}
\date{\today}
	
\maketitle

\pagenumbering{roman}

\section*{Revision History}

\begin{tabularx}{\textwidth}{p{3cm}p{2cm}X}
\toprule {\bf Date} & {\bf Version} & {\bf Notes}\\
\midrule
Date 1 & 1.0 & Notes\\
Date 2 & 1.1 & Notes\\
\bottomrule
\end{tabularx}

~\\
\wss{The intention of the VnV plan is to increase confidence in the software.
However, this does not mean listing every verification and validation technique
that has ever been devised.  The VnV plan should also be a \textbf{feasible}
plan. Execution of the plan should be possible with the time and team available.
If the full plan cannot be completed during the time available, it can either be
modified to ``fake it'', or a better solution is to add a section describing
what work has been completed and what work is still planned for the future.}

\wss{The VnV plan is typically started after the requirements stage, but before
the design stage.  This means that the sections related to unit testing cannot
initially be completed.  The sections will be filled in after the design stage
is complete.  the final version of the VnV plan should have all sections filled
in.}

\newpage

\tableofcontents

\listoftables
\wss{Remove this section if it isn't needed}

\listoffigures
\wss{Remove this section if it isn't needed}

\newpage

\section{Symbols, Abbreviations, and Acronyms}

\renewcommand{\arraystretch}{1.2}
\begin{longtable}{@{}p{0.3\textwidth} p{\dimexpr0.7\textwidth-2\tabcolsep\relax}@{}}
  \toprule
  \textbf{Term} & \textbf{Description}\\
  \midrule
  \endfirsthead
  \toprule
  \textbf{Term} & \textbf{Description}\\
  \midrule
  \endhead
  \bottomrule
  \endfoot
  \bottomrule
  \endlastfoot

  \glsentry{activation}{Activation} & An intermediate value that a model layer computes and passes to the next layer.\\
  \glsentry{ai}{AI} & Artificial intelligence. Software that does tasks that normally need human judgement. In this document it always means an LLM.\\
  \glsentry{api}{API} & Application programming interface. A set of requests one program can make of another. Lake's API is the OpenAI-compatible API on the leader node.\\
  \glsentry{arm}{ARM} & One of the two CPU architectures Lake supports, used by Apple silicon and most mobile devices.\\
  \glsentry{backpressure}{Backpressure} & A signal that a node is receiving tensors faster than it can process them, so the upstream node should slow down.\\
  \glsentry{bottleneck}{Bottleneck} & The node or network connection that takes the longest per token and limits the throughput of the whole Lake network.\\
  \glsentry{contextwindow}{Context window} & The maximum number of tokens, prompt and response together, that a model can handle in one request.\\
  \glsentry{ci}{Continuous integration} & Building and testing the code automatically on every push to the repository.\\
  \glsentry{coordinator}{Coordinator} & The Lake component that forms and monitors the Lake network, serves the API on the leader node, and manages the Inference Engine on each node.\\
  \glsentry{cpu}{CPU} & Central processing unit. A device's general-purpose processor.\\
  \glsentry{deviceowner}{Device owner} & The person whose computer joins a Lake network. They must consent before their device takes part.\\
  \glsentry{dynamicallocation}{Dynamic allocation} & Moving model layers between nodes while the Lake network is running, when a node joins, leaves, or fails.\\
  \glsentry{ffi}{FFI} & Foreign function interface. An interface that lets code in one language call code in another.\\
  \glsentry{flops}{FLOP/s} & Floating-point operations per second. The number of floating-point operations a processor performs per second, used to measure compute performance.\\
  \glsentry{follower}{Follower node} & Any node in a Lake network other than the leader node.\\
  \glsentry{gguf}{GGUF} & The file format used by llama.cpp and ggml to store a model's weights, tokeniser, and metadata. Lake only accepts GGUF models.\\
  \glsentry{gpl}{GPL} & GNU General Public License. The free software licence, version 3.0, that Lake is released under.\\
  \glsentry{gpu}{GPU} & Graphics processing unit. A device's parallel processor, which does the arithmetic of inference.\\
  \glsentry{greedy}{Greedy decoding} & Always choosing the most probable next token, so the same prompt gives the same output every time.\\
  \glsentry{grpc}{gRPC} & A remote procedure call framework from Google. Tensor I/O uses it to send tensors between nodes.\\
  \glsentry{http}{HTTP} & Hypertext Transfer Protocol. The request and response protocol of the web.\\
  \glsentry{inference}{Inference} & Using a model to generate text from a prompt.\\
  \glsentry{inferenceengine}{Inference Engine} & The Lake component on each node that loads the node's partition and runs it to generate output. It is built on ggml.\\
  \glsentry{inferenceuser}{Inference user} & A person, or a tool acting for one, who sends prompts to a Lake network and gets responses back.\\
  \glsentry{ipc}{IPC} & Inter-process communication. Passing data between processes on the same machine.\\
  \glsentry{kvcache}{KV cache} & The key and value tensors each layer stores for past tokens so they are not recomputed. Each node keeps the KV cache for its own layers.\\
  \glsentry{lake}{Lake} & The system this document specifies. It runs one LLM across several devices on a local network using model layer partitioning.\\
  \glsentry{lakeclient}{Lake Client} & The Lake application installed on each device, and the component people use. Every node runs one. The operator forms a Lake network, loads a model, and views performance from the Lake Client on the leader node, and a device owner confirms joining from the Lake Client on their own device.\\
  \glsentry{lakenetwork}{Lake network} & The nodes on one local network that have joined to serve one model.\\
  \glsentry{latency}{Latency} & The time between the start of an inference request and a defined event, such as the first token or the end of the response.\\
  \glsentry{layer}{Layer} & One of the stacked stages of a transformer model. Lake splits a model between nodes at layer boundaries.\\
  \glsentry{leader}{Leader node} & The node whose operator formed the Lake network. It runs the Model Partitioner, serves the API, and monitors the followers.\\
  \glsentry{llm}{LLM} & Large language model. An AI model trained on a large body of text and used to generate text.\\
  \glsentry{partitioning}{Model layer partitioning} & Splitting an LLM's layers into consecutive ranges and giving each range to a different node.\\
  \glsentry{partitioner}{Model Partitioner} & The Lake component on the leader node that decides which layers each node holds, based on each node's memory, compute, and network performance.\\
  \glsentry{moscow}{MoSCoW} & A prioritisation scheme that sorts requirements into Must have, Should have, Could have, and Won't have.\\
  \glsentry{nist}{NIST} & National Institute of Standards and Technology. The United States standards agency.\\
  \glsentry{node}{Node} & A device running the Lake Client that has joined a Lake network.\\
  \glsentry{ocr}{OCR} & Optical character recognition. Converting images of text into machine-readable text.\\
  \glsentry{openaiapi}{OpenAI-compatible API} & An HTTP API with the same request and response formats as the OpenAI API, so clients written for OpenAI work with Lake unchanged.\\
  \glsentry{operator}{Operator} & The person who forms a Lake network, loads a model, and monitors performance from the Lake Client. Also called the Lake operator.\\
  \glsentry{p2p}{P2P} & Peer-to-peer. A direct connection between two nodes that does not go through the leader. Neighbouring nodes send tensors this way.\\
  \glsentry{partition}{Partition} & A consecutive range of a model's layers assigned to one node.\\
  \glsentry{scheme}{Partitioning scheme} & The assignment of partitions to nodes for a whole model.\\
  \glsentry{pegs}{PEGS} & Meyer's four-chapter structure for a requirements document: Project, Environment, Goals, and System.\\
  \glsentry{perplexity}{Perplexity} & A measure of how well a model predicts a body of text, where lower is better. Used to compare a split model with the unsplit one.\\
  \glsentry{pipeda}{PIPEDA} & Personal Information Protection and Electronic Documents Act. Canada's private-sector privacy law.\\
  \glsentry{poc}{PoC} & Proof of concept. The first working demonstration of the system.\\
  \glsentry{prefill}{Prefill} & The first phase of a request, where the whole prompt runs through the model to fill the KV cache before the first output token.\\
  \glsentry{quantization}{Quantisation} & Converting a model's weights to a lower-precision format to reduce memory use.\\
  \glsentry{raci}{RACI} & A chart showing who is Responsible, Accountable, Consulted, and Informed for each piece of work.\\
  \glsentry{rest}{REST} & Representational state transfer. A style of HTTP API where clients read and change named resources.\\
  \glsentry{rpc}{RPC} & Remote procedure call. Calling a function on another machine as if it were local.\\
  \glsentry{rtt}{RTT} & Round-trip time. The time for a message to reach another node and for the reply to come back.\\
  \glsentry{session}{Session} & A period of interaction with Lake in which one or more inference requests are served. For the operator, it runs from forming a Lake network to shutting it down. For an inference user, it is one conversation whose context is kept between prompts.\\
  \glsentry{srs}{SRS} & Software Requirements Specification. This document.\\
  \glsentry{streaming}{Streaming} & Returning a response token by token as it is generated.\\
  \glsentry{ta}{TA} & Teaching assistant. The course staff member who advises the team and assesses its deliverables.\\
  \glsentry{telemetry}{Telemetry} & Usage statistics, crash reports, update checks, or any other data that software sends to its developer or a third party. Lake sends none.\\
  \glsentry{tensor}{Tensor} & A multi-dimensional array of numbers. Weights, activations, and the KV cache are all tensors.\\
  \glsentry{tensorio}{Tensor I/O} & The Lake component that compresses, sends, receives, and decompresses tensors between neighbouring nodes.\\
  \glsentry{token}{Token} & A unit of text processed by an LLM, typically a word or part of a word.\\
  \glsentry{throughput}{Token throughput} & The number of tokens generated per unit time, usually tokens per second.\\
  \glsentry{toolcalling}{Tool calling} & An OpenAI API feature where the model asks the client to run a named function and return the result.\\
  \glsentry{transformer}{Transformer} & The neural network architecture of current LLMs, built from a stack of layers.\\
  \glsentry{ui}{UI} & User interface. The part of a program a person interacts with. Lake's UI is the Lake Client.\\
  \glsentry{vv}{V\&V} & Verification and validation. Verification checks that the system meets its requirements. Validation checks that it is the right system for its stakeholders.\\
  \glsentry{weights}{Weights} & The learned parameters of a model that determine its behaviour. Lake never modifies them.\\
  \glsentry{x86}{x86\_64} & The 64-bit instruction set of Intel and AMD processors. One of the two CPU architectures Lake supports.\\
\end{longtable}

\newpage

\pagenumbering{arabic}

This document ... \wss{provide an introductory blurb and roadmap of the
  Verification and Validation plan}

\section{General Information}

\subsection{Summary}

\wss{Say what software is being tested.  Give its name and a brief overview of
  its general functions.}

\subsection{Objectives}

\wss{State what is intended to be accomplished.  The objective will be around
  the qualities that are most important for your project.  You might have
  something like: ``build confidence in the software correctness,''
  ``demonstrate adequate usability.'' etc.  You won't list all of the qualities,
  just those that are most important.}

\wss{You should also list the objectives that are out of scope.  You don't have 
the resources to do everything, so what will you be leaving out.  For instance, 
if you are not going to verify the quality of usability, state this.  It is also 
worthwhile to justify why the objectives are left out.}

\wss{The objectives are important because they highlight that you are aware of 
limitations in your resources for verification and validation.  You can't do everything, 
so what are you going to prioritize?  As an example, if your system depends on an 
external library, you can explicitly state that you will assume that external library 
has already been verified by its implementation team.}

\subsection{Custom Documents}

\wss{Summarize the CDs (Custom Documents) that were tackled by this project.  CDs 
can include usability testing, code walkthroughs, user documentation, formal
proof, GenderMag personas, Design Thinking, etc.  CDs should have already
been approved by the course instructor as included in your problem statement.
You can use a pull request to update your CDs (in TeamComposition.csv or
Repos.csv) if your plan changes as a result of the VnV planning exercise.}

\subsection{Relevant Documentation}

\wss{Reference relevant documentation.  This will definitely include your SRS
  and your other project documents (design documents, like MG, MIS, etc).  You
  can include these even before they are written, since by the time the project
  is done, they will be written.  You can create BibTeX entries for your
  documents and within those entries include a hyperlink to the documents.}

\citet{SRS}

\wss{Don't just list the other documents.  You should explain why they are relevant and 
how they relate to your VnV efforts.}

\section{Plan}

\wss{Introduce this section.  You can provide a roadmap of the sections to
  come.}

\subsection{Verification and Validation Team}

\wss{Your teammates.  Maybe your supervisor.
  You should do more than list names.  You should say what each person's role is
  for the project's verification.  A table is a good way to summarize this information.}

\subsection{SRS Verification}

\wss{List any approaches you intend to use for SRS verification.  This may
  include ad hoc feedback from reviewers, like your classmates (like your
  primary reviewer), or you may plan for something more rigorous/systematic.}

\wss{If you have a supervisor for the project, you shouldn't just say they will
read over the SRS.  You should explain your structured approach to the review.
Will you have a meeting?  What will you present?  What questions will you ask?
Will you give them instructions for a task-based inspection?  Will you use your
issue tracker?}

\wss{Maybe create an SRS checklist?}

\subsection{Design Verification}

\wss{Plans for design verification}

\wss{The review will include reviews by your classmates}

\wss{Create a checklists?}

\subsection{Verification and Validation Plan Verification}

\wss{The verification and validation plan is an artifact that should also be
verified.  Techniques for this include review and mutation testing.}

\wss{The review will include reviews by your classmates}

\wss{Create a checklists?}

\subsection{Implementation Verification}

\wss{You should at least point to the tests listed in this document and the unit
  testing plan.}

\wss{In this section you would also give any details of any plans for static
  verification of the implementation.  Potential techniques include code
  walkthroughs, code inspection, static analyzers, etc.}

\wss{The final class presentation in CAS 741 could be used as a code
walkthrough.  There is also a possibility of using the final presentation (in
CAS741) for a partial usability survey.}

\subsection{Automated Testing and Verification Tools}

\wss{What tools are you using for automated testing.  Likely a unit testing
  framework and maybe a profiling tool, like ValGrind.  Other possible tools
  include a static analyzer, make, continuous integration tools, test coverage
  tools, etc.  Explain your plans for summarizing code coverage metrics.
  Linters are another important class of tools.  For the programming language
  you select, you should look at the available linters.  There may also be tools
  that verify that coding standards have been respected, like flake9 for
  Python.}

\wss{If you have already done this in the development plan, you can point to
that document.}

\wss{The details of this section will likely evolve as you get closer to the
  implementation.}

\subsection{Software Validation}

\wss{If there is any external data that can be used for validation, you should
  point to it here.  If there are no plans for validation, you should state that
  here.}

\wss{You might want to use review sessions with the stakeholder to check that
the requirements document captures the right requirements.  Maybe task based
inspection?}

\wss{For those capstone teams with an external supervisor, the Rev 0 demo should 
be used as an opportunity to validate the requirements.  You should plan on 
demonstrating your project to your supervisor shortly after the scheduled Rev 0 demo.  
The feedback from your supervisor will be very useful for improving your project.}

\wss{For teams without an external supervisor, user testing can serve the same purpose 
as a Rev 0 demo for the supervisor.}

\wss{This section might reference back to the SRS verification section.}

\section{System Tests}

\wss{There should be text between all headings, even if it is just a roadmap of
the contents of the subsections.}

\subsection{Tests for Functional Requirements}

\wss{Subsets of the tests may be in related, so this section is divided into
  different areas.  If there are no identifiable subsets for the tests, this
  level of document structure can be removed.}

\wss{Include a blurb here to explain why the subsections below
  cover the requirements.  References to the SRS would be good here.}

\subsubsection{Area of Testing1}

\wss{It would be nice to have a blurb here to explain why the subsections below
  cover the requirements.  References to the SRS would be good here.  If a section
  covers tests for input constraints, you should reference the data constraints
  table in the SRS.}
		
\paragraph{Title for Test}

\begin{enumerate}

\item{test-id1\\}

Control: Manual versus Automatic
					
Initial State: 
					
Input: 
					
Output: \wss{The expected result for the given inputs.  Output is not how you
are going to return the results of the test.  The output is the expected
result.}

Test Case Derivation: \wss{Justify the expected value given in the Output field}
					
How test will be performed: 
					
\item{test-id2\\}

Control: Manual versus Automatic
					
Initial State: 
					
Input: 
					
Output: \wss{The expected result for the given inputs}

Test Case Derivation: \wss{Justify the expected value given in the Output field}

How test will be performed: 

\end{enumerate}

\subsubsection{Area of Testing2}

...

\subsection{Tests for Nonfunctional Requirements}

\wss{The nonfunctional requirements for accuracy will likely just reference the
  appropriate functional tests from above.  The test cases should mention
  reporting the relative error for these tests.  Not all projects will
  necessarily have nonfunctional requirements related to accuracy.}

\wss{For some nonfunctional tests, you won't be setting a target threshold for
passing the test, but rather describing the experiment you will do to measure
the quality for different inputs.  For instance, you could measure speed versus
the problem size.  The output of the test isn't pass/fail, but rather a summary
table or graph.}

\wss{Tests related to usability could include conducting a usability test and
  survey.  The survey will be in the Appendix.}

\wss{Static tests, review, inspections, and walkthroughs, will not follow the
format for the tests given below.}

\wss{If you introduce static tests in your plan, you need to provide details.
How will they be done?  In cases like code (or document) walkthroughs, who will
be involved? Be specific.}

\subsubsection{Area of Testing1}
		
\paragraph{Title for Test}

\begin{enumerate}

\item{test-id1\\}

Type: Functional, Dynamic, Manual, Static etc.
					
Initial State: 
					
Input/Condition: 
					
Output/Result: 
					
How test will be performed: 
					
\item{test-id2\\}

Type: Functional, Dynamic, Manual, Static etc.
					
Initial State: 
					
Input: 
					
Output: 
					
How test will be performed: 

\end{enumerate}

\subsubsection{Area of Testing2}

...

\subsection{Traceability Between Test Cases and Requirements}

\wss{Provide a table that shows which test cases are supporting which
  requirements.}

\section{Unit Test Description}

\wss{This section should not be filled in until after the MIS (detailed design
  document) has been completed.}

\wss{Reference your MIS (detailed design document) and explain your overall
philosophy for test case selection.}  

\wss{To save space and time, it may be an option to provide less detail in this section.  
For the unit tests you can potentially layout your testing strategy here.  That is, you 
can explain how tests will be selected for each module.  For instance, your test building 
approach could be test cases for each access program, including one test for normal behaviour 
and as many tests as needed for edge cases.  Rather than create the details of the input 
and output here, you could point to the unit testing code.  For this to work, you code 
needs to be well-documented, with meaningful names for all of the tests.}

\subsection{Unit Testing Scope}

\wss{What modules are outside of the scope.  If there are modules that are
  developed by someone else, then you would say here if you aren't planning on
  verifying them.  There may also be modules that are part of your software, but
  have a lower priority for verification than others.  If this is the case,
  explain your rationale for the ranking of module importance.}

\subsection{Tests for Functional Requirements}

\wss{Most of the verification will be through automated unit testing.  If
  appropriate specific modules can be verified by a non-testing based
  technique.  That can also be documented in this section.}

\subsubsection{Module 1}

\wss{Include a blurb here to explain why the subsections below cover the module.
  References to the MIS would be good.  You will want tests from a black box
  perspective and from a white box perspective.  Explain to the reader how the
  tests were selected.}

\begin{enumerate}

\item{test-id1\\}

Type: \wss{Functional, Dynamic, Manual, Automatic, Static etc. Most will
  be automatic}
					
Initial State: 
					
Input: 
					
Output: \wss{The expected result for the given inputs}

Test Case Derivation: \wss{Justify the expected value given in the Output field}

How test will be performed: 
					
\item{test-id2\\}

Type: \wss{Functional, Dynamic, Manual, Automatic, Static etc. Most will
  be automatic}
					
Initial State: 
					
Input: 
					
Output: \wss{The expected result for the given inputs}

Test Case Derivation: \wss{Justify the expected value given in the Output field}

How test will be performed: 

\item{...\\}
    
\end{enumerate}

\subsubsection{Module 2}

...

\subsection{Tests for Nonfunctional Requirements}

\wss{If there is a module that needs to be independently assessed for
  performance, those test cases can go here.  In some projects, planning for
  nonfunctional tests of units will not be that relevant.}

\wss{These tests may involve collecting performance data from previously
  mentioned functional tests.}

\subsubsection{Module ?}
		
\begin{enumerate}

\item{test-id1\\}

Type: \wss{Functional, Dynamic, Manual, Automatic, Static etc. Most will
  be automatic}
					
Initial State: 
					
Input/Condition: 
					
Output/Result: 
					
How test will be performed: 
					
\item{test-id2\\}

Type: Functional, Dynamic, Manual, Static etc.
					
Initial State: 
					
Input: 
					
Output: 
					
How test will be performed: 

\end{enumerate}

\subsubsection{Module ?}

...

\subsection{Traceability Between Test Cases and Modules}

\wss{Provide evidence that all modules have been considered.}
				
\bibliographystyle{plainnat}

\bibliography{../../refs/References}

\newpage

\section{Appendix}

This is where you can place additional information.

\subsection{Symbolic Parameters}

The definition of the test cases will call for SYMBOLIC\_CONSTANTS.
Their values are defined in this section for easy maintenance.

\subsection{Usability Survey Questions?}

\wss{This is a section that would be appropriate for some projects.}

\newpage{}

\section{Appendix --- Generative AI Use Disclosure}

\wss{ChatGPT (version 5) was used to brainstorm ideas for the template for this disclose section.}

\subsection{AI Tools Used}

\wss{Give the name of the AI tool(s) used and the version.}

\subsection{How and Where Used}

\wss{Explain what the tool(s) were used for.  Examples include: brainstorming,
explaining a technical concept, reviewing the document, generating or modifying
text, generating or modifying diagrams, identifying errors or omissions or
inconsistencies.}

\wss{For each of the uses, identify which parts of the document were affected.}

\wss{You do not need to report every interaction, but you should disclose the 
uses that materially contributed to your work}

\subsection{Verification of AI-Generated Content}

\wss{\begin{itemize}
    \item Describe how the team checked AI-generated or AI-assisted material
    before including it in the document.
    \item In particular, verify factual claims, technical assertions,
    requirements, calculations, references, and other information for which an
    AI system may produce incorrect or fabricated results.
    \item \textbf{The team is responsible for the accuracy, completeness,
consistency, and appropriateness of the submitted document, regardless of
whether material was generated by a human or an AI system.}
\end{itemize}}

\newpage{}

\section{Appendix --- Reflection}

\wss{This section is not required for CAS 741}

The information in this section will be used to evaluate the team members on the
graduate attribute of Lifelong Learning.

\input{../Reflection.text}

\begin{enumerate}
  \item What went well while writing this deliverable? 
  \item What pain points did you experience during this deliverable, and how
    did you resolve them?
  \item What knowledge and skills will the team collectively need to acquire to
  successfully complete the verification and validation of your project?
  Examples of possible knowledge and skills include dynamic testing knowledge,
  static testing knowledge, specific tool usage, Valgrind etc.  You should look to
  identify at least one item for each team member.
  \item For each of the knowledge areas and skills identified in the previous
  question, what are at least two approaches to acquiring the knowledge or
  mastering the skill?  Of the identified approaches, which will each team
  member pursue, and why did they make this choice?
\end{enumerate}

\end{document}
